\section{Solution Method}
\label{sec:solution}

We solve the model by standard log-linearization around the non-stochastic
steady state. The system reduces to three equations in $(\tilde y_t, \pi_t,
i_t)$---the IS curve \eqref{eq:IS_linearized}, the monopsony-adjusted NKPC
\eqref{eq:NKPC_linearized}, and the Taylor rule \eqref{eq:taylor}---
extended with two novel closed-form blocks.

The \textbf{monopsony block} modifies the standard wage-employment
relationship by adding the countercyclical markdown channel. In the
competitive model, the wage Jacobian is $\bar\mu^w \cdot \mathrm{MPL}_n$.
The monopsony block replaces this with:
\begin{equation}
  \frac{\partial w_t}{\partial n_t}\Bigg|_{\text{mono}} =
  \bar\mu^w\,\frac{\partial\mathrm{MPL}_t}{\partial n_t}
  + \bar{\mathrm{MPL}}\cdot
  \underbrace{\frac{\partial\mu^w_t}{\partial u_t}\bigg|_{SS}}_{<\,0}
  \cdot\frac{\partial u_t}{\partial n_t}
  \label{eq:mono_jacobian}
\end{equation}
The second term introduces a negative feedback from employment to wages via
the countercyclical markdown, directly reducing inflationary pressure during
booms. This is what flattens $\kappa^{mono}$ below $\kappa^{comp}$.

The \textbf{fragmentation block} maps shocks to $\phi_t$ into changes in
the full Leontief inverse $\mathbf{L}(\phi_t)$, delivering the TFP
semi-elasticity in closed form (Theorem~\ref{thm:fragmentation}). The
analytical two-sector result (Proposition~\ref{prop:toy}) is the exact
closed form of this block under $S=2$, $R=2$.

Both blocks are analytically tractable and provide the key propositions
without numerical simulation. Quantitative exercises in
Section~\ref{sec:quantitative} use the 10-sector, 5-region calibration;
computational details and convergence checks are in the Supplemental
Appendix.
